%!TEX root = ../QuantumNoiseAndDecoherence.tex
% ──────────────────────────────────────────────────────────────
%  Lecture 5 — Gaussian States and Covariance Matrices
% ──────────────────────────────────────────────────────────────

% Continues Section 4 (Continuous variables).

\subsection{Quadrature operators}\label{sec:quadratures}

So far we parameterised displacement operators using complex numbers
$\vb{\alpha} \in \C^n$.  It is standard in the Gaussian-state
literature to work instead with their real and imaginary parts.

\begin{definition}[Quadrature operators]\label{def:quadratures}
  For a single mode with annihilation operator~$a$, define the
  \textbf{quadrature operators}
  \begin{equation}\label{eq:quadratures}
    \hat{q} = \frac{a + a^\adj}{\sqrt{2}}\,,\qquad
    \hat{p} = \frac{a - a^\adj}{i\sqrt{2}}\,.
  \end{equation}
  Both are hermitian and satisfy $[\hat{q},\, \hat{p}] = i\One$.
\end{definition}

For $n$~modes we collect the $2n$ quadrature operators into a vector
\begin{equation}\label{eq:R-vector}
  \hat{\vb{R}}
  = (\hat{q}_1, \ldots, \hat{q}_n,\;
     \hat{p}_1, \ldots, \hat{p}_n)^\top\,.
\end{equation}
The commutation relations are encoded by the \textbf{symplectic matrix}
\begin{equation}\label{eq:symplectic-matrix}
  [\hat{R}_j,\, \hat{R}_k] = i\,\sigma_{jk}\,\One\,,\qquad
  \sigma = \begin{pmatrix} 0 & \One_n \\ -\One_n & 0 \end{pmatrix}\,.
\end{equation}

\subsubsection{Displacement operator in quadrature notation}

For $\vb{\xi} = (x_1,\ldots,x_n,\, p_1,\ldots,p_n)^\top \in \R^{2n}$
the displacement operator becomes
\begin{equation}\label{eq:displacement-xi}
  D(\vb{\xi})
  = \exp\!\bigl(i\,\vb{\xi}^\top \sigma\, \hat{\vb{R}}\bigr)\,,
\end{equation}
with $\alpha_j = (x_j + i\,p_j)/\sqrt{2}$ connecting to the complex
parameterisation of Lecture~4.  The characteristic function in these
variables is
\begin{equation}\label{eq:char-fn-xi}
  \chi_\rho(\vb{\xi})
  = \tr\bigl(\rho\, D(\vb{\xi})\bigr)\,,\qquad
  \vb{\xi} \in \R^{2n}\,.
\end{equation}

% ──────────────────────────────────────────────────────────────
\subsection{Gaussian states}\label{sec:gaussian-states}

\begin{definition}[Gaussian state]\label{def:gaussian}
  An $n$-mode state~$\rho$ is a \textbf{Gaussian state} if its
  characteristic function is a Gaussian in~$\vb{\xi}$:
  \begin{equation}\label{eq:gaussian-char-fn}
    \eqbox{\chi_\rho(\vb{\xi})
      = \exp\!\Bigl(
          -\tfrac{1}{4}\,\vb{\xi}^\top \gamma\, \vb{\xi}
          - i\,\vb{m}^\top \sigma\, \vb{\xi}
        \Bigr)}\,,
  \end{equation}
  where $\vb{m} \in \R^{2n}$ is the \textbf{displacement vector} and
  $\gamma$ is a real $2n \times 2n$ symmetric matrix called the
  \textbf{covariance matrix}.
\end{definition}

The absence of a constant term is fixed by normalisation:
$\chi_\rho(\vb{0}) = \tr(\rho) = 1$.

The displacement vector records the \textbf{first moments}
\begin{equation}\label{eq:first-moments}
  m_j = \expect{\hat{R}_j} = \tr(\rho\, \hat{R}_j)\,,
\end{equation}
and the covariance matrix records the symmetrised \textbf{second
moments}
\begin{equation}\label{eq:second-moments}
  \gamma_{jk}
  = \expect{\{\hat{R}_j - m_j,\; \hat{R}_k - m_k\}}\,,
\end{equation}
where $\{A,B\} = AB + BA$ is the anticommutator.

\begin{intuition}[Why Gaussian states are remarkable]
  A Gaussian state of $n$~modes lives in an infinite-dimensional
  Hilbert space, yet is completely determined by $\vb{m}$ ($2n$ reals)
  and $\gamma$ ($2n(2n+1)/2$ entries).  All higher moments follow, as
  for a classical Gaussian, making polynomial expectation values
  computable on paper.
\end{intuition}

\begin{remark}
  Superpositions of Gaussian states are \emph{not} generally Gaussian.
  A superposition of the vacuum and a coherent state produces a
  Schr\"odinger-cat-like state, which is non-Gaussian.
\end{remark}

% ──────────────────────────────────────────────────────────────
\subsection{Examples of Gaussian states}\label{sec:gaussian-examples}

\begin{example}[Vacuum state]\label{ex:vacuum-gaussian}
  The $n$-mode vacuum $\rho = \ketbra{\vb{0}}{\vb{0}}$ is Gaussian
  with
  \begin{equation}\label{eq:vacuum-gaussian}
    \vb{m} = \vb{0}\,,\qquad \gamma = \One_{2n}\,.
  \end{equation}
  It is a minimum-uncertainty state:
  $(\Delta q)^2 = (\Delta p)^2 = \tfrac{1}{2}$.
\end{example}

\begin{example}[Coherent states]\label{ex:coherent-gaussian}
  Every coherent state $\ket{\vb{\alpha}}$ is Gaussian.  For a single
  mode with $\alpha = (x_0 + i\,p_0)/\sqrt{2}$:
  \begin{equation}\label{eq:coherent-gaussian}
    \vb{m} = \begin{pmatrix} x_0 \\ p_0 \end{pmatrix}\,,\qquad
    \gamma = \One_2\,.
  \end{equation}
  The covariance matrix equals that of the vacuum; only the
  displacement shifts, consistent with
  $\ket{\alpha} = D(\alpha)\ket{0}$.
\end{example}

% ──────────────────────────────────────────────────────────────
\subsection{Legitimate covariance matrices}\label{sec:legit-cov}

Not every symmetric matrix $\gamma$ corresponds to a physical quantum
state.  Positivity of~$\rho$ imposes a constraint.

\begin{keyresult}[Covariance matrix condition]
  A real symmetric matrix $\gamma \in \R^{2n \times 2n}$ is the
  covariance matrix of a quantum state if and only if
  \begin{equation}\label{eq:cov-condition}
    \eqbox{\gamma + i\sigma \geq 0}\,,
  \end{equation}
  where the inequality is positive semi-definiteness as a hermitian
  matrix.  This is the continuous-variable uncertainty principle.
\end{keyresult}

In particular, $\gamma > 0$ (strictly positive definite): zero
eigenvalues are forbidden.

\begin{proof}[Proof (necessity)]
  For any $\vb{c} \in \C^{2n}$, define
  $A = \sum_j c_j\,(\hat{R}_j - m_j)$.  Since $\rho \geq 0$:
  \begin{equation}\label{eq:cov-proof}
    0 \leq \tr(\rho\, A^\adj A)
    = \sum_{j,k} \bar{c}_j\, c_k\;
      \tr\!\bigl(\rho\,(\hat{R}_j - m_j)(\hat{R}_k - m_k)\bigr)
    = \tfrac{1}{2}\sum_{j,k} \bar{c}_j\, c_k\;
      (\gamma_{jk} + i\sigma_{jk})\,.
  \end{equation}
  Since this holds for all $\vb{c}$, we conclude
  $\gamma + i\sigma \geq 0$.
\end{proof}

\begin{remark}
  The converse also holds: every $(\vb{m}, \gamma)$
  satisfying~\eqref{eq:cov-condition} gives a legitimate Gaussian
  state.  Constructively, for any such~$\gamma$ there exists a system
  of harmonic oscillators whose \emph{thermal state} has that
  covariance matrix, since thermal states of quadratic Hamiltonians
  are always Gaussian.
\end{remark}

% ──────────────────────────────────────────────────────────────
\subsection{Bipartite Gaussian states and reduced states}%
\label{sec:bipartite-gaussian}

Consider two parties, Alice ($A$) and Bob ($B$), each holding
$n$~modes (joint system: $2n$ modes), described by a displacement
vector $\vb{m}_{AB}$ and a $4n \times 4n$ covariance matrix
$\gamma_{AB}$.

\begin{proposition}[Reduced states of Gaussian states]%
\label{prop:reduced-gaussian}
  If $\rho_{AB}$ is Gaussian, then the reduced state
  $\rho_A = \tr_B(\rho_{AB})$ is also Gaussian, with
  \begin{equation}\label{eq:reduced-gaussian}
    \vb{m}_A = (\vb{m}_{AB})_A\,,\qquad
    \gamma_A = (\gamma_{AB})_{AA}\,,
  \end{equation}
  i.e.\ one simply extracts the sub-vector and sub-matrix
  corresponding to Alice's modes.
\end{proposition}

\begin{proof}
  The characteristic function of the reduced state is
  \[
    \chi_{\rho_A}(\vb{\xi}_A)
    = \tr_{AB}\!\bigl(\rho_{AB}\,
        (D(\vb{\xi}_A) \tp \One_B)\bigr)
    = \chi_{\rho_{AB}}(\vb{\xi}_A, \vb{0}_B)\,,
  \]
  using $D(\vb{0}) = \One$.  Setting $\vb{\xi}_B = \vb{0}$ in the
  Gaussian form~\eqref{eq:gaussian-char-fn} yields a Gaussian in
  $\vb{\xi}_A$ with the stated parameters.
\end{proof}

\begin{remark}
  Reducing Gaussian states is trivial (extract a sub-matrix), but the
  reverse---\emph{extending}---is subtler.  Trivial extensions by
  tensoring always exist, but \emph{pure} Gaussian extensions
  (purifications) require a non-trivial construction; they always exist
  and will appear in later lectures.  For tripartite systems $ABC$, it
  is possible that $\rho_{AB}$ and $\rho_{BC}$ are both Gaussian yet
  no global $\rho_{ABC}$ is consistent with both marginals---a
  manifestation of entanglement monogamy.
\end{remark}

% ──────────────────────────────────────────────────────────────
\subsection{Example: the two-mode squeezed state}%
\label{sec:two-mode-squeezed}

The \textbf{two-mode squeezed state} is the canonical entangled
Gaussian state, produced by parametric down-conversion in a nonlinear
crystal.  For two modes $A$, $B$ with
$\hat{\vb{R}} = (\hat{q}_A, \hat{q}_B, \hat{p}_A, \hat{p}_B)^\top$
and squeezing parameter $r \geq 0$:
\begin{equation}\label{eq:tmsv-moments}
  \vb{m} = \vb{0}\,,\qquad
  \gamma =
  \begin{pmatrix}
    \cosh 2r & \sinh 2r & 0 & 0 \\
    \sinh 2r & \cosh 2r & 0 & 0 \\
    0 & 0 & \cosh 2r & -\sinh 2r \\
    0 & 0 & -\sinh 2r & \cosh 2r
  \end{pmatrix}.
\end{equation}
At $r = 0$ this is the two-mode vacuum ($\gamma = \One_4$).
In the Fock basis the corresponding pure state is
\begin{equation}\label{eq:tmsv-fock}
  \ket{\psi_r}
  = \frac{1}{\cosh r}\sum_{n=0}^{\infty}
    (\tanh r)^n\,\ket{n}_A\ket{n}_B\,.
\end{equation}
For large~$r$ the coefficients decay slowly, so $\ket{\psi_r}$
resembles a maximally entangled state up to a photon-number
cutoff---one of the primary entanglement sources in continuous-variable
quantum information.

\begin{exercise}\label{ex:tmsv-reduced}
  Verify Proposition~\ref{prop:reduced-gaussian} for the two-mode
  squeezed state: (a)~extract the sub-matrix of~\eqref{eq:tmsv-moments}
  for mode~$A$; (b)~compute
  $\rho_A = \tr_B(\ketbra{\psi_r}{\psi_r})$ directly
  from~\eqref{eq:tmsv-fock}.  In both cases, confirm
  $\gamma_A = \cosh(2r)\,\One_2$ (a thermal state).
\end{exercise}
